\documentclass{amsart}
% For dealing with references we use the comment environment
\usepackage{verbatim}
\newenvironment{reference}{\comment}{\endcomment}
%\newenvironment{reference}{}{}
\newenvironment{slogan}{\comment}{\endcomment}
\newenvironment{history}{\comment}{\endcomment}

% For commutative diagrams we use Xy-pic
\usepackage[all]{xy}

% We use 2cell for 2-commutative diagrams.
\xyoption{2cell}
\UseAllTwocells

% We use multicol for the list of chapters between chapters
\usepackage{multicol}

% This is generally recommended for better output
\usepackage{lmodern}
\usepackage[T1]{fontenc}

% For cross-file-references

% Package for hypertext links:
\usepackage{hyperref}

% For any local file, say "hello.tex" you want to link to please

% Theorem environments.
%
\theoremstyle{plain}
\newtheorem{theorem}[subsection]{Theorem}
\newtheorem{proposition}[subsection]{Proposition}
\newtheorem{lemma}[subsection]{Lemma}

\theoremstyle{definition}
\newtheorem{definition}[subsection]{Definition}
\newtheorem{example}[subsection]{Example}
\newtheorem{exercise}[subsection]{Exercise}
\newtheorem{situation}[subsection]{Situation}

\theoremstyle{remark}
\newtheorem{remark}[subsection]{Remark}
\newtheorem{remarks}[subsection]{Remarks}

\numberwithin{equation}{subsection}

% Macros
%
\def\lim{\mathop{\mathrm{lim}}\nolimits}
\def\colim{\mathop{\mathrm{colim}}\nolimits}
\def\Spec{\mathop{\mathrm{Spec}}}
\def\Hom{\mathop{\mathrm{Hom}}\nolimits}
\def\Ext{\mathop{\mathrm{Ext}}\nolimits}
\def\SheafHom{\mathop{\mathcal{H}\!\mathit{om}}\nolimits}
\def\SheafExt{\mathop{\mathcal{E}\!\mathit{xt}}\nolimits}
\def\Sch{\mathit{Sch}}
\def\Mor{\mathop{\mathrm{Mor}}\nolimits}
\def\Ob{\mathop{\mathrm{Ob}}\nolimits}
\def\Sh{\mathop{\mathit{Sh}}\nolimits}
\def\NL{\mathop{N\!L}\nolimits}
\def\CH{\mathop{\mathrm{CH}}\nolimits}
\def\proetale{{pro\text{-}\acute{e}tale}}
\def\etale{{\acute{e}tale}}
\def\QCoh{\mathit{QCoh}}
\def\Ker{\mathop{\mathrm{Ker}}}
\def\Im{\mathop{\mathrm{Im}}}
\def\Coker{\mathop{\mathrm{Coker}}}
\def\Coim{\mathop{\mathrm{Coim}}}

% Boxtimes
%
\DeclareMathSymbol{\boxtimes}{\mathbin}{AMSa}{"02}

%
% Macros for moduli stacks/spaces
%
\def\QCohstack{\mathcal{QC}\!\mathit{oh}}
\def\Cohstack{\mathcal{C}\!\mathit{oh}}
\def\Spacesstack{\mathcal{S}\!\mathit{paces}}
\def\Quotfunctor{\mathrm{Quot}}
\def\Hilbfunctor{\mathrm{Hilb}}
\def\Curvesstack{\mathcal{C}\!\mathit{urves}}
\def\Polarizedstack{\mathcal{P}\!\mathit{olarized}}
\def\Complexesstack{\mathcal{C}\!\mathit{omplexes}}
% \Pic is the operator that assigns to X its picard group, usage \Pic(X)
% \Picardstack_{X/B} denotes the Picard stack of X over B
% \Picardfunctor_{X/B} denotes the Picard functor of X over B
\def\Pic{\mathop{\mathrm{Pic}}\nolimits}
\def\Picardstack{\mathcal{P}\!\mathit{ic}}
\def\Picardfunctor{\mathrm{Pic}}
\def\Deformationcategory{\mathcal{D}\!\mathit{ef}}

\setlength{\emergencystretch}{3em}
\begin{document}
\title[Stacks Project, Tag 0FDS]{448 --- Rational-equivalence cycle classes vanish in the adjacent coherent Grothendieck quotient}
\maketitle
\noindent Copyright (C) 2005--2025 Johan de Jong. Stacks Project authors. Source: \href{https://stacks.math.columbia.edu/tag/0FDS}{Tag 0FDS}.

\noindent Editorial difficulty note (not author text). Difficulty H1: A principal divisor must be represented by a difference of coherent quotient sheaves even when its rational function is not globally regular. The proof constructs the ideal of denominators and compares two injections of that ideal, then verifies every codimension-one coefficient through the order/length distance identity..

\noindent Editorial provenance note (not author text). History: excerpt from revision \texttt{a0444\allowbreak 6e57e\allowbreak c1fbc\allowbreak 252a8\allowbreak 71afc\allowbreak ec775\allowbreak 2fb28\allowbreak 07b14}, chow.tex, lines 3933--4017. Reference presentation was adapted; original-author context and core support blocks were retained or added. The mathematical wording is unchanged. Licensed under GNU FDL 1.2 or later; no invariant sections or cover texts. See ../licenses/COPYING. Context blocks reproduce author definitions and supporting results; they are not additional counted problems.

\begin{definition}
\label{algebra-definition-distance}
Let $R$ be a Noetherian local domain of dimension $1$ with
fraction field $K$. Let $V$ be a finite dimensional $K$-vector space.
Let $M$, $M'$ be two lattices in $V$. The {\it distance between
$M$ and $M'$} is the integer
$$
d(M, M') = \text{length}_R(M/M \cap M') - \text{length}_R(M'/M \cap M')
$$
of Lemma \href{https://stacks.math.columbia.edu/tag/02MF}{lemma compare lattices (Tag 02MF)} part (5).
\end{definition}

\begin{definition}
\label{divisors-definition-regular-meromorphic-ideal-denominators}
Let $X$ be a scheme.
Let $\mathcal{L}$ be an invertible $\mathcal{O}_X$-module.
Let $s$ be a regular meromorphic section of $\mathcal{L}$.
The sheaf of ideals $\mathcal{I}$ constructed in
Lemma \href{https://stacks.math.columbia.edu/tag/02P0}{lemma regular meromorphic ideal denominators (Tag 02P0)}
is called the {\it ideal sheaf of denominators of $s$}.
\end{definition}

\begin{situation}
\label{situation-setup}
Here $S$ is a locally Noetherian, and universally catenary scheme.
Moreover, we assume $S$ is endowed with a dimension function
$\delta : S \longrightarrow \mathbf{Z}$.
\end{situation}

\begin{definition}
\label{definition-rational-equivalence}
Let $(S, \delta)$ be as in Situation \ref{situation-setup}.
Let $X$ be a scheme locally of finite type over $S$.
Let $k \in \mathbf{Z}$.
\begin{enumerate}
\item Given any locally finite collection $\{W_j \subset X\}$
of integral closed subschemes with $\dim_\delta(W_j) = k + 1$,
and any $f_j \in R(W_j)^*$ we may consider
$$
\sum (i_j)_*\text{div}(f_j) \in Z_k(X)
$$
where $i_j : W_j \to X$ is the inclusion morphism.
This makes sense as the morphism
$\coprod i_j : \coprod W_j \to X$ is proper.
\item We say that $\alpha \in Z_k(X)$ is {\it rationally equivalent to zero}
if $\alpha$ is a cycle of the form displayed above.
\item We say $\alpha, \beta \in Z_k(X)$ are
{\it rationally equivalent} and we write $\alpha \sim_{rat} \beta$
if $\alpha - \beta$ is rationally equivalent to zero.
\item We define
$$
\CH_k(X) = Z_k(X) / \sim_{rat}
$$
to be the {\it Chow group of $k$-cycles on $X$}. This is sometimes called
the {\it Chow group of $k$-cycles modulo rational equivalence on $X$}.
\end{enumerate}
\end{definition}

\begin{lemma}
\label{lemma-from-chow-to-K}
Let $X$ be a scheme locally of finite type over $(S, \delta)$
as in Situation \ref{situation-setup}. There is a canonical map
$$
\CH_k(X)
\longrightarrow
K_0(\textit{Coh}_{\leq k + 1}(X)/\textit{Coh}_{\leq k - 1}(X))
$$
induced by the map
$Z_k(X) \to K_0(\textit{Coh}_{\leq k}(X)/\textit{Coh}_{\leq k - 1}(X))$
from Lemma \href{https://stacks.math.columbia.edu/tag/02S9}{lemma cycles k group (Tag 02S9)}.
\end{lemma}

\begin{proof}
We have to show that an element $\alpha$ of $Z_k(X)$ which is rationally
equivalent to zero, is mapped to zero in
$K_0(\textit{Coh}_{\leq k + 1}(X)/\textit{Coh}_{\leq k - 1}(X))$.
Write $\alpha = \sum (i_j)_*\text{div}(f_j)$ as in
Definition \ref{definition-rational-equivalence}.
Observe that
$$
\pi = \coprod i_j : W = \coprod W_j \longrightarrow X
$$
is a finite morphism as each $i_j : W_j \to X$ is a closed immersion
and the family of $W_j$ is locally finite in $X$. Hence we may use
Lemma \href{https://stacks.math.columbia.edu/tag/0FDR}{lemma finite cycles k group (Tag 0FDR)} to reduce to the case of $W$.
Since $W$ is a disjoint union of integral scheme, we reduce
to the case discussed in the next paragraph.

\medskip\noindent
Assume $X$ is integral of $\delta$-dimension $k + 1$.
Let $f$ be a nonzero rational function on $X$.
Let $\alpha = \text{div}(f)$. We have to show that
$\alpha$ is mapped to zero in
$K_0(\textit{Coh}_{\leq k + 1}(X)/\textit{Coh}_{\leq k - 1}(X))$.
Let $\mathcal{I} \subset \mathcal{O}_X$ be the ideal of denominators
of $f$, see Divisors, Definition
\ref{divisors-definition-regular-meromorphic-ideal-denominators}.
Then we have short exact sequences
$$
0 \to \mathcal{I} \to \mathcal{O}_X \to \mathcal{O}_X/\mathcal{I} \to 0
$$
and
$$
0 \to \mathcal{I} \xrightarrow{f} \mathcal{O}_X \to
\mathcal{O}_X/f\mathcal{I} \to 0
$$
See Divisors, Lemma
\href{https://stacks.math.columbia.edu/tag/02P0}{lemma regular meromorphic ideal denominators (Tag 02P0)}.
We claim that
$$
[\mathcal{O}_X/\mathcal{I}]_k - [\mathcal{O}_X/f\mathcal{I}]_k =
\text{div}(f)
$$
The claim implies the element $\alpha = \text{div}(f)$ is represented by
$[\mathcal{O}_X/\mathcal{I}] - [\mathcal{O}_X/f\mathcal{I}]$
in $K_0(\textit{Coh}_{\leq k}(X)/\textit{Coh}_{\leq k - 1}(X))$.
Then the short exact sequences show that this element maps to
zero in $K_0(\textit{Coh}_{\leq k + 1}(X)/\textit{Coh}_{\leq k - 1}(X))$.

\medskip\noindent
To prove the claim, let $Z \subset X$ be an integral closed subscheme
of $\delta$-dimension $k$ and let $\xi \in Z$ be its generic point.
Then $I = \mathcal{I}_\xi \subset A = \mathcal{O}_{X, \xi}$
is an ideal such that $fI \subset A$. Now the coefficient of
$[Z]$ in $\text{div}(f)$ is $\text{ord}_A(f)$. (Of course as usual
we identify the function field of $X$ with the fraction field of $A$.)
On the other hand, the coefficient of $[Z]$ in
$[\mathcal{O}_X/\mathcal{I}] - [\mathcal{O}_X/f\mathcal{I}]$
is
$$
\text{length}_A(A/I) - \text{length}_A(A/fI)
$$
Using the distance function of
Algebra, Definition \ref{algebra-definition-distance}
we can rewrite this as
$$
d(A, I) - d(A, fI) = d(I, fI) = \text{ord}_A(f)
$$
The equalities hold by Algebra, Lemmas
\href{https://stacks.math.columbia.edu/tag/02MH}{lemma properties distance function (Tag 02MH)} and
\href{https://stacks.math.columbia.edu/tag/02MI}{lemma order vanishing determinant (Tag 02MI)}.
(Using these lemmas isn't necessary, but convenient.)
\end{proof}
\end{document}
